\documentclass[11pt]{article}
% Shared print layout. The web's provisional design and MathJax are independent.
\usepackage[utf8]{inputenc}
\usepackage[T1]{fontenc}
\usepackage[spanish,es-nodecimaldot]{babel}
\usepackage{lmodern}
\usepackage{amsmath,amssymb}
\usepackage{graphicx}
\usepackage{geometry}
\geometry{a4paper,left=20mm,right=20mm,top=25mm,bottom=22mm,headheight=16pt,headsep=8mm}
\usepackage{fancyhdr,lastpage}
\usepackage{enumitem}
\usepackage{xcolor}
\usepackage{pgfplots}
\pgfplotsset{compat=1.18}
\usepgfplotslibrary{fillbetween}
\usepackage[hidelinks,pdfencoding=auto]{hyperref}
\hypersetup{pdflang=es-ES,pdfauthor={},pdfsubject={Matemáticas Aplicadas a las Ciencias Sociales II}}
\pagestyle{fancy}
\fancyhf{}
\fancyhead[L]{\small Matemáticas Aplicadas a las Ciencias Sociales II}
\fancyhead[R]{\small 2.\textsuperscript{o} Bachillerato}
\fancyfoot[L]{\footnotesize Recuperación y repaso | Edición revisada 1.0}
\fancyfoot[R]{\footnotesize Página \thepage\ de \pageref{LastPage}}
\renewcommand{\headrulewidth}{0.4pt}
\setlength{\parindent}{0pt}
\setlength{\parskip}{0.5em}
\setlist{itemsep=0.3em,topsep=0.3em}
\newcommand{\apartado}[1]{\subsection*{#1}}
\pgfplotsset{
  revision/.style={
    width=0.94\linewidth,height=7.5cm,
    axis lines=middle,grid=major,grid style={gray!18},
    tick label style={font=\small},label style={font=\small},
    xlabel={$x$},ylabel={$y$},samples=160,
    legend style={at={(0.5,-0.16)},anchor=north,draw=none,font=\small},
    clip=true
  }
}

\hypersetup{pdftitle={Recuperación de Inferencia y repaso - Soluciones revisadas}}
\begin{document}

\section*{Ejercicio 1: espera de pacientes}
Soluciones revisadas desde F009. Se conservan los nueve ejercicios y sus puntuaciones originales; falta una regla de selección o conversión de los 22,5 puntos nominales. Se utiliza la normal con \textbf{desviación típica} como segundo parámetro: $N(\mu,\sigma)$. Las varianzas se declaran por separado.

Para la inferencia se asume muestreo aleatorio con observaciones independientes, como requiere el modelo usado en el original. El enunciado de pacientes no describe el procedimiento de selección: esta condición debe comprobarse en una aplicación real.
\begin{enumerate}[label=(\alph*)]
\item La varianza conocida es $\sigma^2=4$, luego $\sigma=2$ minutos. La suma exacta de los 16 datos es
\[
\sum x_i=144,\qquad \bar x=144/16=\boxed{9\text{ minutos}}.
\]
\textbf{Corrección M04:} el original suma 137 y obtiene 8,5625; ambos valores se sustituyen por los de los datos conservados.

Para confianza $c=0,975$, se usa el cuantil positivo $z_c=\Phi^{-1}((1+c)/2)=\Phi^{-1}(0,9875)\approx2,24140273$.
\[
E=z_c\frac{2}{\sqrt{16}}\approx1,12070136,\qquad
IC=\boxed{(7,87929864;\ 10,12070136)}\text{ min}.
\]
Con el cuantil de tabla $z_c\approx2,24$, se obtiene $(7,88;10,12)$ minutos. El procedimiento produce intervalos que cubren la media fija en el 97,5\% de muestras repetidas bajo el modelo; no asigna una probabilidad a esa media después de observar los datos.
\item Para confianza del 90\%, $z_c=\Phi^{-1}(0,95)\approx1,64485363$. Se exige $E\leq0,3$:
\[
n\geq\left(\frac{z_c\,2}{0,3}\right)^2\approx120,24638
\quad\Longrightarrow\quad\boxed{n_{\min}=121}.
\]
Con 120 pacientes el margen es aproximadamente 0,30031 minutos; con 121, 0,29907. Esto confirma la minimalidad. El cuantil de tabla 1,645 lleva al mismo entero.
\end{enumerate}
\clearpage
\section*{Ejercicio 2: precio de las frambuesas}
La varianza poblacional es 9, por lo que $\sigma=3$ euros. No se sustituye por una desviación muestral ni se usa una distribución t.
\begin{enumerate}[label=(\alph*)]
\item Bajo muestreo independiente de la población normal,
\[
\bar X\sim N\left(\mu,\frac{3}{\sqrt{10}}\right),\qquad
\operatorname{Var}(\bar X)=\frac9{10}.
\]
La distribución es normal exacta para cualquier tamaño, no una justificación asintótica por tamaño grande. El original escribe $N(\mu,9/10)$ usando varianza como segundo parámetro; aquí se hace explícita una sola convención para evitar confundir varianza con desviación típica.
\item La suma de los precios conservados es $110$ euros y $\bar x=110/10=\boxed{11}$ euros.

\textbf{Corrección M05:} se reemplazan la suma 109, la media 10,9 y el intervalo desplazado del original.
\[
z_c=\Phi^{-1}(0,985)\approx2,17009038,\quad
E=z_c\frac3{\sqrt{10}}\approx2,05872850.
\]
\[
IC_{97\%}=\boxed{(8,94127150;\ 13,05872850)}\text{ euros}.
\]
Con tabla $z_c\approx2,17$, el margen es aproximadamente 2,05864276 y los extremos son
\[ (8,94135724;13,05864276). \]
A dos decimales, ambos métodos dan $(8,94;13,06)$ euros.
\item Se pide estrictamente $E<1,5$, por lo que
\[
n>\left(\frac{z_c\,3}{1,5}\right)^2\approx18,83717.
\]
El menor entero admisible es $\boxed{19}$. El margen con 18 es aproximadamente 1,53449 euros y con 19, 1,49356 euros. Para una cota estricta se usa $\lfloor q\rfloor+1$, incluso si $q$ fuese entero; no se transforma sin más en una desigualdad no estricta.
\end{enumerate}
\clearpage
\section*{Ejercicio 3: longitud de estrellas de mar}
Los extremos del intervalo se interpretan con la precisión que presenta el enunciado.
\begin{enumerate}[label=(\alph*)]
\item La media muestral es el punto medio:
\[
\bar x=\frac{6,04+8}{2}=\boxed{7,02\text{ cm}}.
\]
\item El margen observado es la semiamplitud:
\[
E=\frac{8-6,04}{2}=\boxed{0,98\text{ cm}}.
\]
Es compatible, al redondear, con el cuantil del 95\%: $\Phi^{-1}(0,975)\,3/6\approx0,979982$ cm.
\item Manteniendo confianza y desviación típica, el margen es proporcional a $1/\sqrt n$. Con el cuantil de tabla 1,96:
\[
E_{49}=1,96\frac3{7}=\boxed{0,84\text{ cm}}.
\]
Sin redondear el cuantil, $E_{49}\approx0,83998456$ cm. También $E_{49}/E_{36}=\sqrt{36/49}=6/7$. No se deduce una nueva media muestral al cambiar de muestra.
\item A tamaño y desviación típica fijos, aumentar la confianza aumenta el cuantil positivo y, por tanto, el margen y la amplitud del intervalo. Se refiere al margen del procedimiento, no al error efectivo desconocido de una estimación concreta.
\end{enumerate}
\clearpage
\section*{Ejercicio 4: llaves sin reemplazamiento}
Se eligen dos llaves distintas entre 90. Hay $\binom{90}{2}=4005$ parejas equiprobables; no son dos ensayos independientes.
\begin{enumerate}[label=(\alph*)]
\item Para que ambas funcionen:
\[
P(2\text{ válidas})=\frac{\binom{60}{2}}{\binom{90}{2}}
=\frac{60}{90}\frac{59}{89}=\boxed{\frac{118}{267}}\approx0,4419.
\]
El original indica 0,440: se corrige el valor decimal conservando su expresión exacta.
\item Por complemento de ninguna válida:
\[
P(\text{alguna válida})=1-\frac{\binom{30}{2}}{\binom{90}{2}}
=1-\frac{29}{267}=\boxed{\frac{238}{267}}\approx0,8914.
\]
No se resta un complemento ya redondeado a 0,108: eso produce el 0,892 del original.
\item Para exactamente una válida hay $60\cdot30$ parejas:
\[
P(1\text{ válida})=\frac{60\cdot30}{4005}
=\boxed{\frac{40}{89}}\approx0,4494.
\]
La cuenta ordenada da lo mismo: $\frac{60}{90}\frac{30}{89}+\frac{30}{90}\frac{60}{89}$. Las probabilidades de cero, una y dos válidas suman 1.
\end{enumerate}
\clearpage
\section*{Ejercicio 5: refrescos y probabilidad condicionada}
Se conserva el año \textbf{2021} como dato histórico del ejercicio; no se sustituye por la fecha actual. ``No caduca en 2021'' significa fuera de ese año, no necesariamente que esté en condiciones de consumo hoy.

Hay 10\,000 latas de cada sabor. Sea $C$ el suceso caducar en 2021.
\[
\begin{array}{c|rr|r}
 &C&C^c&\text{Total}\\\hline
\text{Cola}&1800&8200&10000\\
\text{Naranja}&2400&7600&10000\\
\text{Limón}&3000&7000&10000\\\hline
\text{Total}&7200&22800&30000
\end{array}
\]
\begin{enumerate}[label=(\alph*)]
\item Con elección uniforme entre las latas,
\[
P(C)=\frac{7200}{30000}=\boxed{\frac6{25}=0,24}.
\]
\item Condicionando en las 22\,800 que no caducan en ese año,
\[
P(\text{cola}\mid C^c)=\frac{8200}{22800}
=\boxed{\frac{41}{114}}\approx0,3596.
\]
La selección se hace entre todas las latas; no se inventa otro mecanismo de muestreo.
\end{enumerate}
\clearpage
\section*{Ejercicio 6(a): formular, sin resolver}
Sean $x$ e $y$ las horas de funcionamiento de las cadenas A y B. Se mantiene la exigencia del enunciado de formular \textbf{sin resolver} este modelo:
\[
\boxed{\min Z=1200x+1500y}
\]
\[
\begin{aligned}
10x+7y&\geq400 &&\text{(lavadoras)},\\
5x+6y&\geq280 &&\text{(frigoríficos)},\\
x&\leq2y &&\text{(límite relativo de horas)},\\
x,y&\geq0. &&\text{(horas no negativas)}
\end{aligned}
\]
Las variables son horas continuas. El enunciado no exige enteros ni impone una cota diaria: no se añaden esas condiciones.

Como comprobación de coherencia, no como solución óptima, $(x,y)=(20,30)$ satisface todas las restricciones: produce 410 lavadoras y 280 frigoríficos, y $20\leq60$. Los coeficientes del objetivo son costes por hora y los de las restricciones son producciones por hora.
\clearpage
\section*{Ejercicio 6(b): recinto y mínimo}
\textbf{Solución completada en esta edición:} el original deja este apartado pendiente de representación gráfica.

Las fronteras son $y=(7-x)/2$, $y=4x-1$, $y=2x-4$ e $y=(20-3x)/2$, además de los ejes. La región es el cuadrilátero sombreado, con sus fronteras incluidas.
\begin{center}
\begin{tikzpicture}
\begin{axis}[width=0.86\linewidth,height=7cm,axis lines=middle,xmin=0,xmax=5,ymin=0,ymax=9,tick label style={font=\small},xlabel={$x$},ylabel={$y$},grid=major,legend style={at={(0.5,-0.18)},anchor=north,font=\small,legend columns=2}]
\addplot[fill=blue!15,draw=blue!70!black,thick] coordinates {(1,3)(3,2)(4,4)(2,7)(1,3)};
\addlegendentry{Región factible}
\addplot[domain=0:5,thin] {(7-x)/2};\addlegendentry{$x+2y=7$}
\addplot[domain=0:2.5,thin,dashed] {4*x-1};\addlegendentry{$4x-y=1$}
\addplot[domain=0:5,thin,dotted] {2*x-4};\addlegendentry{$2x-y=4$}
\addplot[domain=0:5,thin,dashdotted] {(20-3*x)/2};\addlegendentry{$3x+2y=20$}
\addplot[only marks,mark=*,blue] coordinates {(1,3)(3,2)(4,4)(2,7)};
\end{axis}
\end{tikzpicture}
\end{center}
\[
\begin{array}{c|c|c}
\text{Vértice}&\text{Fronteras activas}&F=2x+y\\\hline
(1,3)&x+2y=7,\ 4x-y=1&5\\
(3,2)&x+2y=7,\ 2x-y=4&8\\
(4,4)&2x-y=4,\ 3x+2y=20&12\\
(2,7)&4x-y=1,\ 3x+2y=20&11
\end{array}
\]
Se han descartado intersecciones que incumplen alguna desigualdad. Por linealidad sobre este polígono compacto, el mínimo se alcanza en un vértice: $\boxed{F_{\min}=5\text{ en }(1,3)}$. Es único: en las dos aristas incidentes los valores aumentan desde ese vértice.
\clearpage
\section*{Ejercicio 7: parámetros, extremos y primitiva}
\begin{enumerate}[label=(\alph*)]
\item Para pasar por $(2,36)$ y ser estacionaria allí:
\[
8a+2b=32,\qquad12a+b=0
\quad\Longrightarrow\quad\boxed{a=-2,\ b=24}.
\]
Además $f''(2)=6a\cdot2=-24<0$, lo que confirma un máximo relativo; imponer solo $f'(2)=0$ no bastaría en general.
\item Ahora $a=4$, $b=-3$ son datos distintos: $f(x)=4x^3-3x+4$ y $f'(x)=12x^2-3$. Los puntos críticos son $x=\pm1/2$.

La función crece en $(-\infty,-1/2)$ y en $(1/2,\infty)$; decrece en $(-1/2,1/2)$. No se afirma crecimiento sobre la unión como si fuese un único intervalo.
\[
f(-1/2)=-1/2+3/2+4=5,\qquad
f(1/2)=1/2-3/2+4=3.
\]
\textbf{Corrección M24:} máximo relativo $\boxed{(-1/2,5)}$ y mínimo relativo $\boxed{(1/2,3)}$. El original da ordenadas 4 y 2; el cubo conserva el signo de su argumento.
\item Una primitiva es $F(x)=x^4-\frac32x^2+4x+C$. Como $F(2)=18+C=10$, se obtiene $C=-8$:
\[
\boxed{F(x)=x^4-\frac32x^2+4x-8}.
\]
Al derivar se recupera $f$ y al evaluar en 2 se obtiene 10.
\end{enumerate}
\clearpage
\section*{Ejercicio 8: derivadas y área}
\begin{enumerate}[label=(\alph*)]
\item Para $f(x)=(x^2-5)^2e^{3x}$, definida en $\mathbb R$, producto y cadena dan
\[
\boxed{f'(x)=e^{3x}(x^2-5)\bigl(4x+3(x^2-5)\bigr)}.
\]
Para $g(x)=\ln(x^3-5x)/(1-x^2)$, se exige $x^3-5x>0$ y $x\ne\pm1$. Su dominio real es
\[
(-\sqrt5,-1)\cup(-1,0)\cup(\sqrt5,\infty).
\]
Se aplica $(u/v)'=(u'v-uv')/v^2$:
\[
\boxed{g'(x)=\frac{\dfrac{3x^2-5}{x^3-5x}(1-x^2)+2x\ln(x^3-5x)}{(1-x^2)^2}}.
\]
Se corrige la notación intermedia del original, que inserta erróneamente ``$(u/v)'=$'' dentro del numerador; la expresión final del original sí usa correctamente la regla.
\item $h(x)=-x^2+2x+3=-(x+1)(x-3)$ se anula en $-1$ y $3$, y es positiva entre ambas raíces. Por tanto, el área geométrica es
\[
A=\int_{-1}^{3}h(x)\,dx
=\left[-\frac{x^3}{3}+x^2+3x\right]_{-1}^{3}
=9-(-5/3)=\boxed{\frac{32}{3}\text{ unidades}^2}.
\]
No se integran regiones infinitas fuera del intervalo que cierra la parábola con el eje.
\end{enumerate}
\clearpage
\section*{Ejercicio 9: ecuaciones matriciales}
Las matrices son cuadradas, invertibles y del mismo orden, conforme a los productos e hipótesis del enunciado.
\begin{enumerate}[label=(\alph*)]
\item De $B=AX$ y $A=BY$ resulta $A=AXY$. Multiplicando por $A^{-1}$ a la izquierda, $XY=I$. Análogamente $B=BYX$ y $YX=I$. Por tanto,
\[
\boxed{X=Y^{-1}}.
\]
Se usa asociatividad y cancelación lateral, sin intercambiar factores arbitrariamente.
\item Para las matrices dadas, $\det A=1$ y $\det B=-2$. Entonces
\[
A^{-1}=\begin{pmatrix}3&-2\\-1&1\end{pmatrix},\qquad
B^{-1}=\begin{pmatrix}1/2&1/2\\0&-1\end{pmatrix}.
\]
\[
\boxed{X=A^{-1}B=\begin{pmatrix}6&5\\-2&-2\end{pmatrix}},\qquad
\boxed{Y=B^{-1}A=\begin{pmatrix}1&5/2\\-1&-3\end{pmatrix}}.
\]
Comprobaciones directas:
\[
AX=\begin{pmatrix}2&1\\0&-1\end{pmatrix}=B,\qquad
BY=\begin{pmatrix}1&2\\1&3\end{pmatrix}=A,
\]
\[
XY=YX=\begin{pmatrix}1&0\\0&1\end{pmatrix}.
\]
\end{enumerate}
\end{document}
